\documentclass[9pt]{extarticle}
\usepackage[margin=0.5in]{geometry}
\usepackage{booktabs}
\usepackage{enumitem}
\usepackage{xcolor}
\usepackage{titlesec}
\usepackage{helvet}
\renewcommand{\familydefault}{\sfdefault}
\definecolor{navy}{HTML}{203764}
\titleformat{\section}{\normalfont\bfseries\color{navy}}{\thesection}{0.5em}{}[\vspace{-0.6em}]
\titlespacing*{\section}{0pt}{5pt}{2pt}
\setlist{nosep, leftmargin=1.2em, topsep=1pt, partopsep=0pt}
\pagestyle{empty}
\linespread{0.94}
\begin{document}

\begin{center}
{\LARGE\bfseries\color{navy} Adobe Inc. (NASDAQ: ADBE) --- Equity Pitch}\\[2pt]
{\large\bfseries Recommendation: LONG \quad Target price: \$500 \quad Price: \$247.09 (close, Oct 19, 2026) \quad Implied upside: +102\%}\\[2pt]
Archit Chawla \textbar\ Market cap \$96.2B \textbar\ EV \$96.9B \textbar\ Shares out 389.2M (Sep 18, 2026) \textbar\ Street: Hold, avg target \$272--284
\end{center}
\vspace{-4pt}
\hrule
\vspace{2pt}

\section*{Variant perception}
The market prices Adobe as an AI loser whose creative franchise melts away: at \$247 the stock is worth almost exactly my disruption scenario (\$242), in which revenue growth fades to low-single digits and the FCF margin compresses toward 30\%. The filings show the opposite trajectory --- growth re-accelerating toward 13\%, a $\sim$41\% FCF margin, and a buyback retiring roughly 6\% of the share count every year.

\section*{Business snapshot}
Adobe is the incumbent in digital content creation and experience: Digital Media (Creative Cloud, Document Cloud/Acrobat) sells subscriptions to creators, knowledge workers and enterprises; Digital Experience sells the marketing/analytics stack to the same base. $\sim$90\% gross margin, subscription-led, with \$7.09B of current deferred revenue at Q3 FY2026 --- revenue already contracted, not yet recognized. No dividend (verified in the filings); all capital return is buybacks; essentially no net debt (\$0.72B).

\section*{Thesis}
\begin{enumerate}
\item \textbf{Growth is re-accelerating, not rolling over.} Revenue grew +12.0\%, +12.7\% and +12.9\% YoY in Q1--Q3 FY2026, up from +10.5\% in FY2025; LTM revenue is \$26.0B. Deferred revenue grew +11.1\% YoY to \$7.09B --- the forward book compounds at double digits, which is not what a disrupted franchise looks like.
\item \textbf{The cash machine is intact.} LTM free cash flow is \$10.59B, a 40.8\% margin; FY2025 operating cash flow was \$10.03B on capex of just 0.8\% of revenue. AI investment runs through opex --- it is already inside these numbers, not a capex wall still to come.
\item \textbf{The buyback is a second return engine.} \$11.28B repurchased in FY2025 (114\% of FCF); \$35.7B over five years; shares outstanding down 13.9\% since Jan 2024 (452.0M $\rightarrow$ 389.2M). In the model, diluted shares fall 427M $\rightarrow$ 344M by FY2028E, so GAAP EPS compounds \$18.97 $\rightarrow$ \$29.67 while net income grows a more modest 6--19\% a year. (SBC of $\sim$\$2.1B/yr is charged through the share-count roll --- a real cost, not added back.)
\item \textbf{The multiple prices a funeral.} Today: 9.7$\times$ EV/LTM EBITDA and 9.1$\times$ price/LTM FCF. Adobe's own fiscal-year-end multiples, FY2021--FY2025 (EDGAR fundamentals, calendar year-end market caps): 15.7--40.6$\times$ EV/EBITDA and 15.2--39.1$\times$ P/FCF; medians 25.6$\times$ and 25.0$\times$. Even FY2025's de-rated 15.2$\times$ P/FCF on FY2027E FCF/share gives \$534.
\end{enumerate}

\section*{Valuation}
DCF on FY2027E--FY2031E FCF from a linked three-statement model (revenue \$26.71B FY26E, \$29.65B FY27E, \$32.77B FY28E --- an 11.3\% FY25--28E CAGR with growth \emph{faded}; FCF \$11.24B / \$12.91B / \$14.37B). WACC 11.0\%: risk-free 5.1\%, beta 1.42 (Finnhub/MarketBeat, Oct 20, 2026), ERP 4.5\% $\rightarrow$ cost of equity 11.5\%; after-tax cost of debt 3.3\%; terminal growth 3.0\%.
\begin{center}\vspace{2pt}
\begin{tabular}{@{}llll@{}}
\toprule
\textbf{Scenario} & \textbf{Assumptions} & \textbf{Value/share} & \textbf{vs \$247.09} \\
\midrule
Bull & FCF +8\%, WACC 10.5\%, g 3.5\% & \$609 & +147\% \\
\textbf{Base (target)} & \textbf{Model FCF, WACC 11.0\%, g 3.0\%} & \textbf{\$500} & \textbf{+102\%} \\
Bear & FCF $-$18\%, WACC 11.5\%, g 2.5\% & \$369 & +49\% \\
AI disruption & Growth $\rightarrow$ LSD, FCF margin $\rightarrow$ $\sim$30\%, WACC 11.5\%, g 1.5\% & \$242 & $-$2\% \\
\bottomrule
\end{tabular}
\end{center}
\vspace{1pt}
Cross-checks cluster with the DCF: 20.0$\times$ FY2027E GAAP EPS (\$24.36) = \$487; FY2025's own P/FCF on FY2027E FCF/share = \$534. Sensitivity (WACC $\pm$100bp, g 2.5--3.5\%): \$425--609. Consensus (\$272--284) sits between my disruption and bear cases --- it is paying for partial disruption. The Q4 FY2026 stub FCF ($\sim$\$3.8B, $\approx$\$10/share) is excluded from the DCF, a deliberate conservatism.

\section*{Catalysts}
\textbf{Dec 9, 2026:} Q4 FY2026 results + FY2027 guidance (Q4 guide: non-GAAP EPS \$6.30--6.35; my model's FY26 non-GAAP EPS proxy is \$24.64 vs the \$24.45--24.50 guide). \textbf{Buyback pace} in that print ($\sim$\$2.6B assumed for Q4): at \$247, every \$1B retires $\sim$1\% of the company. \textbf{AI monetization disclosures} (Firefly volumes, Acrobat AI attach) across FY2027 prints --- the evidence the market says it is waiting for.

\section*{Risks --- and what would change my call}
\begin{itemize}
\item \textbf{AI-native competition} (Canva, Figma, generative image/video) taking Creative seats. A real risk --- and already the price.
\item \textbf{Seat-growth slowdown.} Tripwires: deferred revenue growth below +5\% YoY for two consecutive quarters, or revenue growth below +9\% --- pillar 1 breaks; I would exit to the disruption valuation.
\item \textbf{Guidance credibility.} A cut to FY2027 guidance on Dec 9 breaks the thesis's hinge.
\item \textbf{SBC creep.} Stock comp is 8.2\% of revenue (FY2025); toward 10\%, the buyback runs harder just to stand still.
\item \textbf{Volatility.} Beta is 1.42; size for a revisit of the \$190 June 2026 low.
\end{itemize}

\section*{Methodology \& verification}
All historicals are SEC EDGAR company facts (CIK 0000796343), pulled Oct 20, 2026; periods keyed on end dates (EDGAR's \texttt{fy} field is the filing's year). Adobe's fiscal year ends the Friday nearest Nov 30 (FY2025 ended Nov 28, 2025). The Excel model (live formulas) was recalculated in LibreOffice and verified: FY2025 revenue (\$23,769M), net income (\$7,130M) and operating cash flow (\$10,031M) tie exactly to EDGAR; the balance sheet balances (check $=$ 0) in every forecast year; balance-sheet cash equals the cash-flow roll; the DCF target equals the sensitivity grid's center cell (\$499.91). FY2026E $=$ 9M actuals + a Q4 stub grown off derived Q4 FY2025, each rate an assumption cell with a rationale. The IS amortization line is derived as OperatingExpenses $-$ R\&D $-$ S\&M $-$ G\&A (\$157M, FY2025) --- the identically named EDGAR tag is a cash-flow concept that breaks the tie-out by $\sim$\$150M.

\section*{Limitations}
Beta is used as reported (1.42, itself inflated by the de-rating) --- conservative for the target. Multiples history uses calendar year-end market caps (companiesmarketcap.com) vs late-November fiscal year-ends, and year-end snapshots rather than averages. FY2026E embeds one estimated quarter. The model is consolidated (EDGAR facts do not support a segment build). SBC is charged through dilution rather than as a cash cost in FCF --- the standard treatment, stated openly.

\vspace{4pt}
\hrule
\vspace{2pt}
{\footnotesize Independent student research project, prepared for educational purposes. Not investment advice. Figures should be verified against the underlying SEC filings before any investment decision. Full model, data pulls and sources: github.com/ArchitChawla26/Adobe-Equity-Pitch}

\end{document}
